\section{Conclusão}
O objetivo geral foi atendido pela comparação, em seis receptores humanos associados à pele, entre o ligante de menor FEB do conjunto comercial e coocristalizado e o composto para prospecção selecionado pela sobreposição dos contatos. Os pares selecionados foram L2U e glicose-6-sulfato em MMP-9; GW475151 e glicose-6-sulfato em MMP-1; L2U e galactose-6-sulfato em MMP-2; batimastat e ácido ferúlico em TYRP1; ARP-100 e Dha em elastase; e cromoglicato e ácido gálico em hialuronidase 1.

As duplas compartilharam de três a cinco resíduos de contato. A maior sobreposição ocorreu entre L2U e glicose-6-sulfato em MMP-9, com cinco resíduos em comum. A glicose-6-sulfato também compartilhou quatro resíduos com GW475151 em MMP-1. Em MMP-2, L2U e galactose-6-sulfato tiveram quatro resíduos em comum; nos demais receptores, cada dupla compartilhou três. Em todos os pares, o ligante comparador apresentou FEB mais negativa. Em 2PE4, o ácido gálico foi selecionado pela sobreposição de contatos com cromoglicato, embora Dha tenha apresentado a menor FEB entre os compostos para prospecção.

Os objetivos específicos foram cumpridos: foram delimitados 17 compostos para prospecção com base em publicações sobre SPE e materiais microalgais relacionados; foram selecionadas estruturas cristalográficas de seis receptores; e foram realizadas as simulações com a função Vina para elastase e hialuronidase 1 e AD4 com mapas AD4Zn para MMP-1, MMP-2, MMP-9 e TYRP1. A análise das pontuações e dos relatórios PLIP, apoiada pela inspeção no PyMOL, permitiu identificar e comparar os resíduos envolvidos nas poses selecionadas.

Os resultados caracterizam os perfis de docking dos compostos isolados e indicam quais podem ser priorizados em ensaios bioquímicos. A confirmação de efeitos sobre a atividade enzimática e a caracterização da SPE de \textit{T. chuii} dependem de etapas experimentais.
